\section{Double Pendulum}

In this problem a planar double pendulum is considered: a point mass $m_1$ hangs from a fixed pivot (the origin) on a massless rigid rod of length $l_1$, and a second point mass $m_2$ hangs from the first on a massless rigid rod of length $l_2$. The angles $\theta_1$ and $\theta_2$ are measured from the downward vertical, and gravity $g$ acts along $-y$. The system is solved in two ways, as an unconstrained problem in generalized coordinates and as a constrained problem in Cartesian coordinates, and the two solutions are compared.

\subsection{Euler--Lagrange Formulation}

The two degrees of freedom are $q = (\theta_1, \theta_2)$. The coordinates of the masses are

\begin{align}
    (x_1, y_1) &= \left( l_1\sin\theta_1,\; -l_1\cos\theta_1 \right) \\
    (x_2, y_2) &= \left( l_1\sin\theta_1 + l_2\sin\theta_2,\; -l_1\cos\theta_1 - l_2\cos\theta_2 \right).
\end{align}

The total kinetic energy of the system is $\mathcal{T} = \mathcal{T}_{m_1} + \mathcal{T}_{m_2}$ where

\begin{align}
    \mathcal{T}_{m_1} &= \frac{1}{2}m_1 l_1^2\dot\theta_1^2 \\
    \mathcal{T}_{m_2} &= \frac{1}{2}m_2\left( l_1^2\dot\theta_1^2 + l_2^2\dot\theta_2^2 + 2l_1l_2\dot\theta_1\dot\theta_2\cos(\theta_1-\theta_2) \right)
\end{align}

and the total potential energy is

\begin{equation}
    \mathcal{V} = -(m_1+m_2)\,g\,l_1\cos\theta_1 - m_2\,g\,l_2\cos\theta_2.
\end{equation}

The Lagrangian $\mathcal{L} = \mathcal{T} - \mathcal{V}$ is therefore

\begin{align}
    \mathcal{L} &= \frac{1}{2}(m_1+m_2)l_1^2\dot\theta_1^2 + \frac{1}{2}m_2l_2^2\dot\theta_2^2 + m_2l_1l_2\dot\theta_1\dot\theta_2\cos(\theta_1-\theta_2) \nonumber \\
&\quad + (m_1+m_2)\,g\,l_1\cos\theta_1 + m_2\,g\,l_2\cos\theta_2
\end{align}

and the Euler--Lagrange equation is asserted for each generalized coordinate $q_i \in \{\theta_1, \theta_2\}$:

\begin{equation}
    \frac{d}{dt}\left( \frac{\partial \mathcal{L}}{\partial \dot{q}_i} \right) - \frac{\partial \mathcal{L}}{\partial q_i} = 0.
\end{equation}

The partial derivatives with respect to $\theta_1$ and $\dot\theta_1$ are

\begin{align}
    \frac{\partial \mathcal{L}}{\partial \theta_1} &= -m_2l_1l_2\dot\theta_1\dot\theta_2\sin(\theta_1-\theta_2) - (m_1+m_2)\,g\,l_1\sin\theta_1 \\
    \frac{\partial \mathcal{L}}{\partial \dot\theta_1} &= (m_1+m_2)l_1^2\dot\theta_1 + m_2l_1l_2\dot\theta_2\cos(\theta_1-\theta_2)
\end{align}

and with respect to $\theta_2$ and $\dot\theta_2$

\begin{align}
    \frac{\partial \mathcal{L}}{\partial \theta_2} &= m_2l_1l_2\dot\theta_1\dot\theta_2\sin(\theta_1-\theta_2) - m_2\,g\,l_2\sin\theta_2 \\
    \frac{\partial \mathcal{L}}{\partial \dot\theta_2} &= m_2l_2^2\dot\theta_2 + m_2l_1l_2\dot\theta_1\cos(\theta_1-\theta_2).
\end{align}

Expressions for $\ddot{\theta}_1$ and $\ddot{\theta}_2$ as explicit functions of time were gotten using Julia's symbolic solver library \texttt{Symbolics}. Since this is the first time a symbolic solver is introduced in this report, these results were (painstakingly) confirmed analytically as well.

\begin{minted}{julia}
using Symbolics

# Define variables
@variables t
D = Differential(t)

@variables m1 m2 l1 l2 g

# Define functions of time
@variables θ1(t) θ2(t)

# Define derivatives
θ1_dot = D(θ1)
θ2_dot = D(θ2)
θ1_ddot = D(θ1_dot)
θ2_ddot = D(θ2_dot)

# Define static variables to evaluate partial derivatives
@variables q1 q2 q1_dot q2_dot

# Define static varibles to eval partial derivatives
L_static = 0.5 * (m1 + m2) * l1^2 * q1_dot^2 + 
           0.5 * m2 * l2^2 * q2_dot^2 + 
           m2 * l1 * l2 * q1_dot * q2_dot * cos(q1 - q2) + 
           (m1 + m2) * g * l1 * cos(q1) + 
           m2 * g * l2 * cos(q2)

# Evaulate partial differentials
dL_dq1_dot = expand_derivatives(Differential(q1_dot)(L_static))
dL_dq2_dot = expand_derivatives(Differential(q2_dot)(L_static))
dL_dq1 = expand_derivatives(Differential(q1)(L_static))
dL_dq2 = expand_derivatives(Differential(q2)(L_static))

# Substitute the time-dependent dynamic functions into the expressions
sub_back = Dict(
    q1 => θ1,
    q2 => θ2,
    q1_dot => θ1_dot,
    q2_dot => θ2_dot
)

pL_pθ1_dot = substitute(dL_dq1_dot, sub_back)
pL_pθ2_dot = substitute(dL_dq2_dot, sub_back)
pL_pθ1 = substitute(dL_dq1, sub_back)
pL_pθ2 = substitute(dL_dq2, sub_back)

# Form the Euler-Lagrange Equation
eq1 = expand_derivatives(D(pL_pθ1_dot)) - pL_pθ1 ~ 0
eq2 = expand_derivatives(D(pL_pθ2_dot)) - pL_pθ2 ~ 0

# Solve
solutions = Symbolics.solve_for(
    [eq1, eq2], [θ1_ddot, θ2_ddot]
)

# Simplify the final explicit equations of motion
θ1_ddot_explicit = simplify(solutions[1])
θ2_ddot_explicit = simplify(solutions[2])    
\end{minted}

Taking the total time derivatives of $\partial\mathcal{L}/\partial\dot\theta_i$, the cross terms proportional to $\dot\theta_1\dot\theta_2\sin(\theta_1-\theta_2)$ cancel in part, as they did for the single pendulum. Dividing the first equation by $l_1$ and the second by $m_2l_2$ leaves two coupled second-order equations, which are written in matrix form with $\Delta = \theta_1-\theta_2$:


The determinant of the matrix is $l_1l_2\left(m_1 + m_2\sin^2\Delta\right) > 0$, so the system can be inverted in closed form (Cramer's rule) at every instant. It is solved with the Tsitouras 5/4 Runge--Kutta method (\texttt{Tsit5}) in Julia, with the state $u = (\theta_1, \theta_2, \dot\theta_1, \dot\theta_2)$.

\begin{align}
\ddot{\theta}_1 &= \frac{g m_1 \sin\theta_1 + g m_2 \sin\theta_1 - g m_2 \sin\theta_2 \cos(\theta_1 - \theta_2) + \dot{\theta}_2^2 l_2 m_2 \sin(\theta_1 - \theta_2) + \dot{\theta}_1^2 l_1 m_2 \sin(\theta_1 - \theta_2) \cos(\theta_1 - \theta_2)}{-l_1 m_1 - l_1 m_2 + l_1 m_2 \cos^2(\theta_1 - \theta_2)} \\
\ddot{\theta}_2 &= \frac{g m_1 \sin\theta_2 + g m_2 \sin\theta_2 - g m_1 \sin\theta_1 \cos(\theta_1 - \theta_2) - g m_2 \sin\theta_1 \cos(\theta_1 - \theta_2) - \dot{\theta}_1^2 l_1 m_1 \sin(\theta_1 - \theta_2) - \dot{\theta}_1^2 l_1 m_2 \sin(\theta_1 - \theta_2) - \dot{\theta}_2^2 l_2 m_2 \sin(\theta_1 - \theta_2) \cos(\theta_1 - \theta_2)}{-l_2 m_1 - l_2 m_2 + l_2 m_2 \cos^2(\theta_1 - \theta_2)}
\end{align}

\begin{minted}{julia}

Base.@kwdef struct DoublePendulumParams
    m1::Float64 = 1.0
    m2::Float64 = 1.0
    l1::Float64 = 1.0
    l2::Float64 = 1.0
    g::Float64  = 9.81
end

# Euler-Lagrange double pendulum: A(θ) θ̈ = b(θ, ω)
function double_pendulum!(du, u, p, t)
    θ1, θ2, ω1, ω2 = u
    m1, m2, l1, l2, g = p.m1, p.m2, p.l1, p.l2, p.g
    Δ = θ1 - θ2

    s, c = sincos(Δ) # A little Julia quirk

    a11 = (m1 + m2) * l1
    a12 = m2 * l2 * c
    a21 = l1 * c
    a22 = l2
    b1 = -m2 * l2 * ω2^2 * s - (m1 + m2) * g * sin(θ1)
    b2 =  l1 * ω1^2 * s      - g * sin(θ2)

    det = a11 * a22 - a12 * a21      *# = l1 l2 (m1 + m2 sin²Δ)*

    du[1] = ω1
    du[2] = ω2
    du[3] = (b1 * a22 - a12 * b2) / det
    du[4] = (a11 * b2 - a21 * b1) / det
return nothing
end
\end{minted}

\subsection{Differential-Algebraic Formulation}

Alternatively, the double pendulum can be described with the Cartesian coordinates of both masses, $\mathbf{q} = (x_1, y_1, x_2, y_2)^\top$. This introduces two redundant coordinates, so the rigid rods must be imposed as holonomic constraints

\begin{align}
    \Phi_1(\mathbf{q}) &= x_1^2 + y_1^2 - l_1^2 = 0 \\
    \Phi_2(\mathbf{q}) &= (x_2-x_1)^2 + (y_2-y_1)^2 - l_2^2 = 0.
\end{align}

The rods exert constraint forces on the masses. Writing $\Delta x = x_2 - x_1$, $\Delta y = y_2 - y_1$ and introducing the Lagrange multipliers $\lambda_1, \lambda_2$, Newton's second law for each mass is

\begin{align}
    m_1\ddot{x}_1 &= -2\lambda_1 x_1 + 2\lambda_2\Delta x \\
    m_1\ddot{y}_1 &= -m_1 g - 2\lambda_1 y_1 + 2\lambda_2\Delta y \\
    m_2\ddot{x}_2 &= -2\lambda_2\Delta x \\
    m_2\ddot{y}_2 &= -m_2 g - 2\lambda_2\Delta y
\end{align}

We have four Newton equations for six unknowns: four accelerations and two multipliers. Differentiating the rod constraints gives the two extra equations.

The first derivatives are the velocity constraints:
\begin{align}
\dot\Phi_1 &= 2(x_1\dot x_1+y_1\dot y_1)=0,\\
\dot\Phi_2 &= 2(\Delta x\,\Delta\dot x+\Delta y\,\Delta\dot y)=0.
\end{align}
Here $\Delta\dot x=\dot x_2-\dot x_1$, and similarly for $y$.

The second derivatives are the acceleration constraints:
\begin{align}
\ddot\Phi_1 &= 2(x_1\ddot x_1+y_1\ddot y_1)
              +2(\dot x_1^2+\dot y_1^2)=0,\\
\ddot\Phi_2 &= 2(\Delta x\,\Delta\ddot x+\Delta y\,\Delta\ddot y)
              +2(\Delta\dot x^2+\Delta\dot y^2)=0.
\end{align}
Here $\Delta\ddot x=\ddot x_2-\ddot x_1$, and similarly for $y$.

This DAE system  may be written as a single matrix equation. Over here, however, it is expressed as two equations here for the sake of brevity.

\begin{align}
    \begin{pmatrix}
        1 & 0 & 0 & 0 \\
        0 & 1 & 0 & 0 \\
        0 & 0 & 1 & 0 \\
        0 & 0 & 0 & 1
    \end{pmatrix}
    \begin{pmatrix}
        \dot{x}_1 \\
        \dot{y}_1 \\
        \dot{x}_2 \\
        \dot{y}_2 \\
    \end{pmatrix}
    &=
    \begin{pmatrix}
        v_{x1} \\
        v_{y1} \\
        v_{x2} \\
        v_{y2} \\
    \end{pmatrix} \\
    \begin{pmatrix}
    m_1 & 0 & 0 & 0 & 2x_1 & -2\Delta x \\
    0 & m_1 & 0 & 0 & 2y_1 & -2\Delta y \\
    0 & 0 & m_2 & 0 & 0 & 2\Delta x \\
    0 & 0 & 0 & m_2 & 0 & 2\Delta y \\
    x_1 & y_1 & 0 & 0 & 0 & 0 \\
    -\Delta x & -\Delta y & \Delta x & \Delta y & 0 & 0
    \end{pmatrix}
    \begin{pmatrix}
        \dot{v}_{x1} \\
        \dot{v}_{y1} \\
        \dot{v}_{x2} \\
        \dot{v}_{y2} \\
        \lambda_1 \\
        \lambda_2 \\
    \end{pmatrix}
    &=
    \begin{pmatrix}
        0 \\
        -m_1g \\
        0 \\
        -m_2g \\
        -\left(v_{x1}^2 + v_{y1}^2 \right) \\
        -\left(\Delta v_x^2 + \Delta v_y^2 \right)
    \end{pmatrix}
\end{align}



\begin{figure}[htbp]
    \centering
    % First Subfigure
    \begin{subfigure}[b]{0.48\textwidth}
        \centering
        \includegraphics[width=\textwidth]{Problems/DoublePendu/Lagrangian_Formalism/figures/end_effector_trajectory.png} % Replace with your image file
        \caption{Lagrangian Mechanics}
        \label{fig:sub_first}
    \end{subfigure}
    \hfill % Adds horizontal spacing between the subfigures
    % Second Subfigure
    \begin{subfigure}[b]{0.48\textwidth}
        \centering
        \includegraphics[width=\textwidth]{Problems/DoublePendu/DAEs/figures/trajectory.png} % Replace with your image file
        \caption{Differential Algebraic Equations}
        \label{fig:sub_second}
    \end{subfigure}
    
    \caption{Solving the double pendulum with different methods: $l_1 = l_2 = m_1 = m_2 = 1.0$ and $\theta_1^\text{initial} = \theta_2^\text{initial} = \frac{\pi}{2}$, $\omega_1^\text{initial} = \omega_2^\text{initial} = 0.0$}
    \label{fig:main_label}
\end{figure}

\subsection{Newton-Euler Formalism}
Must be done, and will be done!